% Compiles on its own (pulls ../main.tex's preamble via subfiles) or as part of the book.
\documentclass[../main.tex]{subfiles}
\begin{document}

\chapter{The Task, and the Data That Defines It}

This chapter is the theory gate (R2) for the first thing we build: the data pipeline.
Before writing a loader, we settle what the loader must produce and why every one of its
choices is load-bearing. In multispectral detection, more results have been quietly ruined
by data-handling mistakes than by bad architectures. We start where the mistakes start.

\section{What we are building, precisely}

We have two cameras observing one scene: a visible-light (RGB) camera and a long-wave
infrared (LWIR, ``thermal'') camera. Each emits an image. We want a single neural network
that consumes \emph{both} and outputs a list of boxes around every pedestrian. The research
question is not ``can we detect pedestrians'' --- that is solved --- but \emph{where inside
the network the two streams should meet, and what operation should merge them}, and
\emph{how that choice behaves when the two images stop lining up}.

The benchmark is KAIST (Hwang et al., CVPR 2015): $95{,}328$ RGB--thermal image pairs
captured from a vehicle through a beam splitter, at $640\times512$, $20$\,Hz, day and night.

\section{The label file, decoded}

KAIST ships one XML file per frame. A single annotated object looks like this:

\begin{lstlisting}
<object>
  <name>person</name>
  <bndbox> <x>562</x> <y>219</y> <w>43</w> <h>91</h> </bndbox>
  <occlusion>0</occlusion>
</object>
\end{lstlisting}

Three facts here decide whether our numbers mean anything.

\paragraph{1. The box format is \texttt{(x, y, w, h)}.} That is the \emph{top-left corner}
plus \emph{width} and \emph{height} --- not two corners. Nearly every detection framework
speaks $(x_1,y_1,x_2,y_2)$ internally, so the very first thing our parser does is convert:
\[
  x_2 = x + w, \qquad y_2 = y + h.
\]
Skip this and every box is the right size but anchored to the wrong place --- a bug that
still ``trains,'' just to a worse model, which is the hardest kind to notice.

\paragraph{2. Only \texttt{person} is a positive.} The labels contain four names:
\texttt{person}, \texttt{cyclist}, \texttt{people} (a crowd too dense to box individually),
and \texttt{person?} (annotator uncertain). The KAIST protocol scores \emph{only}
\texttt{person}. The other three are \emph{ignore regions}: a detection that lands on one is
thrown away before scoring --- neither counted as a hit nor charged as a false alarm.

\begin{lesson}
Ignore regions must be \emph{kept}, not deleted. Their whole job is to absorb detections
that would otherwise be scored as false positives. A loader that filters the annotations
down to \texttt{person} and discards the rest will inflate the false-positive count and
report a worse miss rate than the identical model evaluated correctly. The information that
protects your score lives in the labels you were tempted to throw away.
\end{lesson}

\paragraph{3. Occlusion is a tag, not a guess.} The \texttt{<occlusion>} field is
$0$ (none), $1$ (partial), or $2$ (heavy). The headline ``Reasonable'' subset keeps
pedestrians with occlusion $0$ or $1$ and height $\geq 55$ pixels, and ignores the rest.
We store the tag now so the evaluator can apply the subset later; we do not filter at load
time, because training and evaluation want different subsets of the same data.

\section{The two images}

Both modalities arrive as $512\times640$, 8-bit JPEGs. The thermal image is stored with
three channels, but it originates from a single-sensor grayscale stream; the minute
per-channel differences are JPEG chroma artifacts, not real color. So thermal is
\emph{effectively single-channel}, and how we turn one channel into the three a
conventional backbone expects is a real modeling decision we defer to a later chapter --- but
we flag it here because it is exactly the kind of ``obvious'' preprocessing step where a
silent train/test mismatch hides.

\begin{lesson}
Trust the decoded array, not the documentation. Much of the KAIST literature states the
frames are $640\times480$; the actual JPEGs and the XML \texttt{<size>} field both say
$640\times512$. A resize built on the wrong assumed size distorts every box coordinate by a
few percent --- again, invisibly, because the model still trains.
\end{lesson}

\section{What the loader must return}

For one frame index, the loader returns a paired sample:
\[
  \big(\, I_{\text{rgb}} \in \mathbb{R}^{3\times H\times W},\;\;
          I_{\text{thm}} \in \mathbb{R}^{3\times H\times W},\;\;
          B \in \mathbb{R}^{N\times4},\;\;
          c \in \{0,1\}^{N},\;\;
          o \in \{0,1,2\}^{N} \,\big),
\]
where $B$ are boxes in $(x_1,y_1,x_2,y_2)$ pixel coordinates, $c$ marks each box positive
(\texttt{person}) or ignore, and $o$ is the occlusion tag. Because KAIST is beam-splitter
registered, \emph{one} box set serves both modalities --- for now. The moment we start
injecting misalignment (Chapter~[misalignment]), that assumption is exactly what we perturb,
which is why the loader keeps the two images as separate tensors even though they share a
label: the separation is the seam along which the whole headline experiment later opens.

\section{The first slice to prove}

The thinnest end-to-end slice for a detection project is one sample, drawn and looked at:
image in $\rightarrow$ boxes parsed $\rightarrow$ boxes drawn on \emph{both} modalities
$\rightarrow$ a human confirms they land on people. That is the build for this chapter, and it
passes: the boxes land on pedestrians in both streams.

The original draft of this section set a second, more ambitious exit check --- \emph{find five
pairs where the RGB and thermal boxes visibly fail to coincide}, to see the residual
misalignment by eye. Building it taught the opposite lesson, and it is worth recording exactly
because the first instinct was wrong.

\begin{lesson}
The residual RGB--thermal misalignment on KAIST (a few pixels on a $\sim$90\,px body,
$3$--$5\%$) is below what the naked eye reliably resolves --- and on a \emph{night} sequence
the point collapses entirely, because RGB can barely see the pedestrian at all, so there is no
shared object whose edges you could compare between the two modalities. We confirmed this with
a false-colour overlay (thermal in red, RGB in cyan): the pedestrians came out pure red
(thermal-only) and the painted crosswalk pure cyan (RGB-only). That image is a vivid lesson in
\emph{modality complementarity} --- each sensor sees what the other cannot --- but it is not a
misalignment measurement, and pretending otherwise would be exactly the kind of
see-what-you-expect error the whole project is built to avoid.
\end{lesson}

The honest conclusion reframes the misalignment axis and, in doing so, justifies the project's
central design choice. Misalignment on this benchmark cannot be \emph{eyeballed}; it must be
\emph{measured} --- either against per-modality boxes (Zhang et al.'s KAIST-Paired annotations,
which exist precisely because the shift is real but small) or, as we will do, by \emph{injecting}
a known shift and measuring how each fusion operator degrades (the Phase~4 headline). We did not
choose ``inject and measure'' over ``inspect'' for convenience; we chose it because inspection
demonstrably does not work here. The revised exit check for this chapter is therefore the one we
actually passed: the loader is correct, the boxes land on people in both modalities, and the
false-colour overlay makes the complementarity of the two sensors --- the reason fusion exists
at all --- plainly visible.

\end{document}
